\documentclass[11pt]{article}
\usepackage[margin=1in]{geometry}
\usepackage{amsmath,amssymb,amsthm,booktabs,longtable,array,calc}
\usepackage[hidelinks]{hyperref}
\usepackage[T1]{fontenc}
\usepackage[utf8]{inputenc}
\newtheorem{theorem}{Theorem}[section]
\newtheorem{proposition}[theorem]{Proposition}
\newtheorem{lemma}[theorem]{Lemma}
\newtheorem{corollary}[theorem]{Corollary}
\theoremstyle{definition}
\newtheorem{definition}[theorem]{Definition}
\theoremstyle{remark}
\newtheorem*{remark}{Remark}
\providecommand{\tightlist}{\setlength{\itemsep}{0pt}\setlength{\parskip}{0pt}}
\begin{document}
\appendix
\section{\texorpdfstring{Technical Closure for
\(MMMMVV\)}{Technical Closure for MMMMVV}}\label{appendix-d.-technical-closure-for-mmmmvv}

This appendix supplies the detailed necessity, existence, full-validity,
and uniqueness bridges used in Theorem 7.1.

\subsection{D.1. Row-normal-form column words and horizontal
actions}\label{d.1.-row-normal-form-column-words-and-horizontal-actions}

Put \(o_r=2r-1\), \(e_r=2r\) for \(1\le r\le k\). The plus-normal-form
column word is \[
W_{2s-1}=(o_s,o_{s-1},\ldots,o_1,e_1,\ldots,e_{s-1},o_{s+1},\ldots,o_k,e_k,\ldots,e_{s+1},e_s),
\] \[
W_{2s}=(o_{s+1},\ldots,o_k,e_k,\ldots,e_{s+1},o_s,\ldots,o_1,e_1,\ldots,e_s).
\] For \(MMMMVV\), horizontal crease \(H_{o_r}\) has opener \(A_{o_r}\)
and closer \(B_{o_r}\), while \(H_{e_r}\) has opener \(B_{e_r}\) and
closer \(A_{e_r}\).

\subsection{D.2. Necessity: the four-edge
cycle}\label{d.2.-necessity-the-four-edge-cycle}

Define \(X_p(i,j)\) by \(e_j\) preceding \(o_i\) in \(W_p\). Direct
inspection gives \[
X_{2s-1}(i,j)\iff (i>s\ \text{and}\ j<s),
\] \[
X_{2s}(i,j)\iff (i\le s\ \text{and}\ j>s).
\] If \(1<a<m\) and \(a\ne b\), choose \[
(i,j)=(s+1,s-1)\quad\text{when }a=2s-1,
\] \[
(i,j)=(s,s+1)\quad\text{when }a=2s.
\] The legal ranges are \(2\le s\le k-1\) in the odd case and
\(1\le s\le k-1\) in the even case. The following complete case table
verifies \(X_a=1\) and \(X_b=0\):

\begin{longtable}[]{@{}
  >{\raggedright\arraybackslash}p{(\columnwidth - 6\tabcolsep) * \real{0.2500}}
  >{\raggedright\arraybackslash}p{(\columnwidth - 6\tabcolsep) * \real{0.2500}}
  >{\raggedright\arraybackslash}p{(\columnwidth - 6\tabcolsep) * \real{0.2500}}
  >{\raggedright\arraybackslash}p{(\columnwidth - 6\tabcolsep) * \real{0.2500}}@{}}
\toprule\noalign{}
\begin{minipage}[b]{\linewidth}\raggedright
parity of \((a,b)\)
\end{minipage} & \begin{minipage}[b]{\linewidth}\raggedright
condition
\end{minipage} & \begin{minipage}[b]{\linewidth}\raggedright
witness
\end{minipage} & \begin{minipage}[b]{\linewidth}\raggedright
reason \(X_b=0\)
\end{minipage} \\
\midrule\noalign{}
\endhead
\bottomrule\noalign{}
\endlastfoot
odd/odd & \(a<b\) & \((s+1,s-1)\) & \(X_b\) would require \(i>t\), but
\(t\ge s+1=i\) \\
odd/odd & \(a>b\) & \((s+1,s-1)\) & \(X_b\) would require \(j<t\), but
\(t\le s-1=j\) \\
odd/even & any \(b\ne a\) & \((s+1,s-1)\) & if \(b=2t\), then \(X_b=1\)
would require \(s+1\le t\) and \(s-1>t\), simultaneously \\
even/odd & any \(b\ne a\) & \((s,s+1)\) & if \(b=2t-1\), then \(X_b=1\)
would require \(s>t\) and \(s+1<t\), simultaneously \\
even/even & \(a<b\) & \((s,s+1)\) & \(X_b\) would require \(j>t\), but
\(t\ge s+1=j\) \\
even/even & \(a>b\) & \((s,s+1)\) & \(X_b\) would require \(i\le t\),
but \(t\le s-1<i=s\) \\
\end{longtable}

Set \(o=2i-1\) and \(e=2j\). In any merge with pivots \((a,b)\), \[
A_e\prec A_o
\] comes from \(X_a=1\) in the upper row. Since \(o\) is odd and \(e\)
even, horizontal orientation gives \[
A_o\prec B_o,\qquad B_e\prec A_e.
\] Since \(X_b=0\), the strict lower row gives \(B_o\prec B_e\). Thus \[
A_e\prec A_o\prec B_o\prec B_e\prec A_e,
\] impossible in a total order. At \(k=2\), the only interior upper
pivot is \(a=2\), the witness is \((i,j)=(1,2)\), and the same cycle is
\(A_4\prec A_1\prec B_1\prec B_4\prec A_4\).

\subsection{D.3. Five explicit existence
families}\label{d.3.-five-explicit-existence-families}

For a column \(c\), define \[
Q_c=(A_c,B_c)\quad(c\text{ odd}),\qquad
Q_c=(B_c,A_c)\quad(c\text{ even}).
\] Empty ranges are omitted.

If \(W_p=(c_1,\ldots,c_{2k})\), define the diagonal construction \[
D(k,p)=Q_{c_1}Q_{c_2}\cdots Q_{c_{2k}},
\] which has pivots \((p,p)\).

For \(1\le s\le k\) define \[
\begin{aligned}
L_{\rm odd}(k,s)=&\ A_{o_1}\cdots A_{o_s}\mid B_{o_s}\cdots B_{o_1}\mid B_{e_1}\cdots B_{e_{s-1}}\\
&\mid Q_{o_{s+1}}\cdots Q_{o_k}\mid Q_{e_k}\cdots Q_{e_s}\mid A_{e_{s-1}}\cdots A_{e_1},
\end{aligned}
\] with pivots \((1,2s-1)\), and for \(1\le s\le k-1\) define \[
\begin{aligned}
L_{\rm even}(k,s)=&\ A_{o_1}\cdots A_{o_{s+1}}\mid B_{o_{s+1}}\mid Q_{o_{s+2}}\cdots Q_{o_k}\\
&\mid Q_{e_k}\cdots Q_{e_{s+1}}\mid B_{o_s}\cdots B_{o_1}\mid B_{e_1}\cdots B_{e_s}\mid A_{e_s}\cdots A_{e_1},
\end{aligned}
\] with pivots \((1,2s)\).

For \(1\le s\le k\) define \[
\begin{aligned}
R_{\rm odd}(k,s)=&\ A_{o_k}\cdots A_{o_s}\mid B_{o_s}\mid Q_{o_{s-1}}\cdots Q_{o_1}\mid Q_{e_1}\cdots Q_{e_{s-1}}\\
&\mid B_{o_{s+1}}\cdots B_{o_k}\mid B_{e_k}\cdots B_{e_s}\mid A_{e_s}\cdots A_{e_k},
\end{aligned}
\] with pivots \((m,2s-1)\), and for \(1\le s\le k-1\) define \[
\begin{aligned}
R_{\rm even}(k,s)=&\ A_{o_k}\cdots A_{o_{s+1}}\mid B_{o_{s+1}}\cdots B_{o_k}\mid B_{e_k}\cdots B_{e_{s+1}}\\
&\mid Q_{o_s}\cdots Q_{o_1}\mid Q_{e_1}\cdots Q_{e_s}\mid A_{e_{s+1}}\cdots A_{e_k},
\end{aligned}
\] with pivots \((m,2s)\).

The index ranges partition the four label classes, so every displayed
word is a strict total order of all \(4k\) labels. Deleting all \(B\)
labels or all \(A\) labels gives exactly the claimed row-normal-form
words. The true formula overlaps are \((1,1)\) and \((m,m)\), where the
diagonal and boundary words are literally identical. The pairs \((1,m)\)
and \((m,1)\) are ordinary boundary members of \(L_{\rm odd}(k,k)\) and
\(R_{\rm odd}(k,1)\). All formulas remain literal at \(k=2\) under the
empty-range convention.

\subsection{D.4. Horizontal stack validity of the
constructions}\label{d.4.-horizontal-stack-validity-of-the-constructions}

For \(MMMMVV\), the horizontal actions are

\[
A_{o_r}=\operatorname{push}(o_r),\qquad B_{o_r}=\operatorname{pop}(o_r),\qquad
B_{e_r}=\operatorname{push}(e_r),\qquad A_{e_r}=\operatorname{pop}(e_r).
\]

Each \(Q_c\) is therefore an immediate push/pop pair above the
pre-existing stack. The five construction families can be checked
literally.

\begin{itemize}
\tightlist
\item
  \textbf{Diagonal \(D(k,p)\).} Every \(Q_{c_j}\) pushes and immediately
  pops \(c_j\), so the stack is empty before and after every block.
\item
  \textbf{Left odd \(L_{\rm odd}(k,s)\).} The initial
  \(A_{o_1},\ldots,A_{o_s}\) pushes the odd core \((o_1,\ldots,o_s)\),
  and \(B_{o_s},\ldots,B_{o_1}\) pops it. Then
  \(B_{e_1},\ldots,B_{e_{s-1}}\) pushes the even core
  \((e_1,\ldots,e_{s-1})\). All later \(Q\) blocks act above that core,
  and \(A_{e_{s-1}},\ldots,A_{e_1}\) pops it.
\item
  \textbf{Left even \(L_{\rm even}(k,s)\).} The first odd run pushes
  \((o_1,\ldots,o_{s+1})\); \(B_{o_{s+1}}\) pops its top, leaving
  \((o_1,\ldots,o_s)\). The intervening odd and even \(Q\) blocks act
  above that core. Then \(B_{o_s},\ldots,B_{o_1}\) empties the odd core,
  \(B_{e_1},\ldots,B_{e_s}\) pushes a fresh even core, and
  \(A_{e_s},\ldots,A_{e_1}\) pops it.
\item
  \textbf{Right odd \(R_{\rm odd}(k,s)\).} The first run pushes
  \((o_k,o_{k-1},\ldots,o_s)\); \(B_{o_s}\) removes the top. The \(Q\)
  blocks act above the residual odd core \((o_k,\ldots,o_{s+1})\), which
  is then popped by \(B_{o_{s+1}},\ldots,B_{o_k}\). Finally
  \(B_{e_k},\ldots,B_{e_s}\) pushes the even core \((e_k,\ldots,e_s)\)
  and \(A_{e_s},\ldots,A_{e_k}\) pops it.
\item
  \textbf{Right even \(R_{\rm even}(k,s)\).} The initial odd core
  \((o_k,\ldots,o_{s+1})\) is pushed and then popped completely. The run
  \(B_{e_k},\ldots,B_{e_{s+1}}\) pushes an even core; all \(Q\) blocks
  act above it; the final run \(A_{e_{s+1}},\ldots,A_{e_k}\) pops it.
\end{itemize}

All empty endpoint ranges are literal, so the same verification covers
\(s=1\), \(s=k\), and \(s=k-1\) where those parameters are legal. In
every family there is no underflow, every closer matches the current
top, and the final stack is empty. Consequently every horizontal
directed relation holds and every two horizontal intervals are disjoint
or nested, never strictly alternating.

\subsection{\texorpdfstring{D.5. Remaining \(U/L\) full-validity
obligations}{D.5. Remaining U/L full-validity obligations}}\label{d.5.-remaining-ul-full-validity-obligations}

Proposition 6.1 already handles all directed vertical relations and
same-row equal-target pairs. H-stack validity handles the horizontal
system. Therefore only $U_i/L_j$ with equal seam parity remain.

For the diagonal family, each column is a contiguous block $Q_c$ in
the row-normal-form column order $W_p$. If $i\ne j$ have equal parity,
then the adjacent-column pairs $\{i,i+1\}$ and $\{j,j+1\}$ are disjoint.
Strict alternation of the actual $U/L$ endpoints would force alternation
of those column pairs in $W_p$, contradicting row-local validity from
Proposition 6.1. For $i=j$, the two adjacent column blocks occur with
opposite internal orientations, and direct inspection shows one vertical
interval nested in the other.

For a boundary family fix the upper $A$ chain. Let $\rho(c)$ be the
position of $A_c$, and put
\[
\sigma(c)=\#\{A_d:A_d\prec B_c\}.
\]

\begin{lemma}[Slots and cross-row alternation]\label{lem:per2-slots}
Suppose an upper seam has endpoint ranks $\alpha<\beta$ in the upper
$A$ chain and a lower seam has endpoint slots $x\le y$. The two seams
strictly alternate if and only if exactly one of $x,y$ belongs to the
integer interval $[\alpha,\beta-1]$. In particular, alternation is
impossible if $x=y$, or if $x+y=\alpha+\beta-1$.
\end{lemma}
\begin{proof}
A lower face lies strictly between the upper endpoints exactly when at
least $\alpha$, but fewer than $\beta$, upper faces precede it, i.e. its
slot belongs to $[\alpha,\beta-1]$. Two intervals with four distinct
endpoints alternate exactly when precisely one endpoint of the second
interval lies inside the first. This gives the criterion and the same-slot
assertion. If $x+y=\alpha+\beta-1$ and $x\ge\alpha$, then
$y\le\beta-1$, so both slots lie inside. If $x<\alpha$, integrality gives
$x\le\alpha-1$ and hence $y\ge\beta$, so the lower endpoints lie on the
two exterior sides. Neither case gives exactly one interior endpoint.
\end{proof}

For upper chain $W_1$,
\[
\rho_1(o_r)=r,\qquad \rho_1(e_r)=2k-r+1,
\]
so
\[
S_r^o=[r,2k-r]\quad(1\le r\le k),\qquad
S_r^e=[r+1,2k-r]\quad(1\le r\le k-1).
\]
For upper chain $W_m$,
\[
\rho_m(o_r)=k-r+1,\qquad \rho_m(e_r)=k+r,
\]
so
\[
T_r^o=[k-r+1,k+r-1]\quad(1\le r\le k),\qquad
T_r^e=[k-r,k+r-1]\quad(1\le r\le k-1).
\]
The sums of the two endpoints in these four shell families are,
respectively, $2k$, $2k+1$, $2k$, and $2k-1$.

The required \(\sigma\) values are as follows; the three columns
correspond to \(t<s\), \(t=s\), and \(t>s\).

\begin{longtable}[]{@{}llrrr@{}}
\toprule\noalign{}
construction & column & \(t<s\) & \(t=s\) & \(t>s\) \\
\midrule\noalign{}
\endhead
\bottomrule\noalign{}
\endlastfoot
\(L_{\rm odd}(k,s)\) & \(\sigma(o_t)\) & \(s\) & \(s\) & \(t\) \\
& \(\sigma(e_t)\) & \(s\) & \(2k-s\) & \(2k-t\) \\
\(L_{\rm even}(k,s)\) & \(\sigma(o_t)\) & \(2k-s\) & \(2k-s\) & \(t\) \\
& \(\sigma(e_t)\) & \(2k-s\) & \(2k-s\) & \(2k-t\) \\
\(R_{\rm odd}(k,s)\) & \(\sigma(o_t)\) & \(k-t+1\) & \(k-s+1\) &
\(k+s-1\) \\
& \(\sigma(e_t)\) & \(k+t-1\) & \(k+s-1\) & \(k+s-1\) \\
\(R_{\rm even}(k,s)\) & \(\sigma(o_t)\) & \(k-t+1\) & \(k-s+1\) &
\(k-s\) \\
& \(\sigma(e_t)\) & \(k+t-1\) & \(k+s-1\) & \(k-s\) \\
\end{longtable}

These values are obtained directly by counting preceding \(A\) labels in
the displayed construction words. They give the following complete
lower-seam classification.

\begin{itemize}
\tightlist
\item
  In \(L_{\rm odd}\), odd seam \(2t-1\) is a point for \(t<s\), is
  \(S_s^o\) for \(t=s\), and is \(S_t^o\) for \(t>s\); even seam \(2t\)
  is a point for \(t<s\), is \(S_s^e\) for \(t=s\), and is \(S_t^e\) for
  \(t>s\).
\item
  In \(L_{\rm even}\), odd seam \(2t-1\) is a point for \(t\le s\) and
  is \(S_t^o\) for \(t>s\); even seam \(2t\) is a point for \(t<s\), is
  \(S_s^e\) for \(t=s\), and is \(S_t^e\) for \(t>s\).
\item
  In \(R_{\rm odd}\), odd seam \(2t-1\) is \(T_t^o\) for \(t\le s\) and
  a point for \(t>s\); even seam \(2t\) is \(T_t^e\) for \(t<s\) and a
  point for \(t\ge s\).
\item
  In \(R_{\rm even}\), odd seam \(2t-1\) is \(T_t^o\) for \(t\le s\) and
  a point for \(t>s\); even seam \(2t\) is \(T_t^e\) for \(t\le s\) and
  a point for \(t>s\).
\end{itemize}

Thus every non-point lower seam is exactly a member of the corresponding
upper shell family. For an arbitrary upper seam of the same parity, its
shell and the two lower endpoint slots belong to the same one of the four
families above and therefore have the same endpoint sum. If the upper
shell is $[\alpha,\beta-1]$, the lower slots satisfy
$x+y=\alpha+\beta-1$, so Lemma~\ref{lem:per2-slots} excludes strict
alternation. Point seams are excluded by the same lemma. This proves all
remaining $U_i/L_j$ obligations for arbitrary separation of the seam
indices. The argument uses the constant endpoint sum; numerical
containment of slot intervals by itself is not used.

\subsection{D.6. H-local legality is not
H-completability}\label{d.6.-h-local-legality-is-not-h-completability}

For fixed row chains, represent a partial H merge by consumed row-prefix
lengths and the stack \(S\) of open columns. A next row head is
H-locally legal if it obeys one push/pop step. It is H-completable only
if the child state admits a full continuation ending with every face
consumed and the stack empty.

These notions differ. At \(k=2\) with \((a,b)=(1,1)\), after prefix
\(A_1\) both \(A_3\) and \(B_1\) are locally legal, but only \(B_1\)
admits completion. Therefore the uniqueness proof never uses the false
rule ``every locally legal next event is globally forced.''

\begin{lemma}[Diagonal fixed-row uniqueness]\label{lem:per2-diagonal}
For the canonical assignment $\texttt{MMMMVV}$ and every
$1\le p\le m$, the two row orders with pivots $(p,p)$ have exactly one
complete H-valid merge, namely $D(k,p)$.
\end{lemma}
\begin{proof}
Write their common column order as $W_p=(c_1,\ldots,c_n)$.
Assume inductively that the first $r-1$ columns have been processed as
$Q_{c_1}\cdots Q_{c_{r-1}}$. The horizontal stack is empty and both row
heads have column $c_r$. The opener must be emitted first: it is
$A_{c_r}$ for odd $c_r$, and $B_{c_r}$ for even $c_r$.

The other row is now waiting at its matching closer. Suppose the row
that emitted the opener advances instead to its next column $d$.
If that next event is a closer, its matching opener has not appeared,
since it occurs after $c_r$ in the blocked opposite row. If it is an
opener, its label $d$ is pushed above $c_r$. In any complete LIFO trace,
$d$ would then have to close before $c_r$, whereas the opposite row
requires the closer of $c_r$ before the closer of $d$. Both alternatives
are impossible. If this row is exhausted, there is no alternative event.
Hence the matching closer follows immediately, producing $Q_{c_r}$,
emptying the stack, and restoring the induction hypothesis.
The resulting word is $D(k,p)$, which is H-valid by its immediate
push--pop blocks. The proof includes $p=1,m$ as well as interior pivots.
\end{proof}

\subsection{D.7. Projected determinism and first
disagreement}\label{d.7.-projected-determinism-and-first-disagreement}

In any complete H-valid merge, deleting all events of one column parity
leaves a legal stack trace. For the odd projection, the upper row fixes
the opener stream and the lower row fixes the closer stream. For the
even projection these roles are exchanged. A fixed opener stream and a
fixed closer stream have at most one complete legal stack trace:
if the next requested closer has not opened, the opener stream is forced
through that label; once it is the top label, closing it is necessary,
since another push would bury the fixed next closer. A requested closer
that is already buried cannot be completed. Thus both parity projections
are uniquely determined whenever they exist.

For a boundary pivot pair, let $T$ be the explicit H-valid construction
from D.3--D.4. Suppose a different complete H-valid merge $T'$ exists,
and consider their first divergence. Their common prefix is a prefix
of the known word $T$. An upper-odd/lower-odd divergence contradicts
uniqueness of the odd projection, and an upper-even/lower-even divergence
contradicts uniqueness of the even projection. An upper-even/lower-odd
divergence would require two distinct closers to be simultaneously legal
at one stack top. The only remaining possibility is an upper-odd opener
versus a lower-even opener.

\subsection{D.8. Opener windows on the explicit
constructions}\label{d.8.-corrected-cross-parity-windows}

The following table lists all prefixes of the four explicit boundary
constructions at which both row heads are openers. In the stack column,
entries are written from bottom to top.
\[
\begin{array}{c|c|c|c}
T&\text{range}&\text{stack}&\text{row heads}\\\hline
L_{\rm odd}(k,s)&2\le s\le k-1,\ 1\le t\le s-1
 &(e_1,\ldots,e_{t-1})&(A_{o_{s+1}},B_{e_t})\\
L_{\rm even}(k,s)&\text{none}&-&-\\
R_{\rm odd}(k,s)&\text{none}&-&-\\
R_{\rm even}(k,s)&1\le s\le k-1,\ s+1\le t\le k
 &(e_k,e_{k-1},\ldots,e_{t+1})&(A_{o_s},B_{e_t}).
\end{array}
\]
In every listed window, $T$ chooses the lower-even opener.
For $L_{\rm odd}$, the initial odd core is opened and closed before the
lower even core $e_1,\ldots,e_{s-1}$ is opened. During this last phase
the next upper face is $A_{o_{s+1}}$ exactly when $s<k$. The range is
empty for $s=1$, and for $s=k$ the upper head is already an even closer.
For $R_{\rm even}$, the initial odd core is likewise emptied; the lower
even core is then opened from $e_k$ down to $e_{s+1}$ while the upper
head remains $A_{o_s}$, giving the second displayed range.
For $L_{\rm even}$ and $R_{\rm odd}$, all upper odd openers have appeared
before the first lower even opener is reached. Immediate $Q_c$ blocks
and the final closing phases create no further opener/opener windows.
This exhausts prefixes of the explicit constructions, not arbitrary
H-locally legal dead-end prefixes. That is exactly the scope required
by the first-divergence argument.

\subsection{D.9. Wrong-switch barriers and completion of
uniqueness}\label{d.9.-wrong-switch-barriers}

At a left-odd window, suppose $T'$ chooses $A_{o_{s+1}}$ instead of
$B_{e_t}$. Before the lower row can reach $B_{o_{s+1}}$, it must open
$e_t,\ldots,e_{s-1}$ above the newly opened odd label. Consequently
LIFO requires $A_{e_{s-1}}\prec B_{o_{s+1}}$. The fixed row orders and
horizontal precedence then give the impossible cycle
\[
 B_{o_{s+1}}
 \mathrel{\prec_B} B_{e_k}
 \mathrel{\prec_H} A_{e_k}
 \mathrel{\prec_A} A_{e_{s-1}}
 \mathrel{\prec_{\rm LIFO}} B_{o_{s+1}}.
\]
Here $\prec_A,\prec_B$ are the prescribed upper and lower row orders;
$\prec_H$ is horizontal opener-before-closer precedence.

At a right-even window, choosing $A_{o_s}$ forces the later lower opener
$e_{s+1}$ above it before $B_{o_s}$ can appear. LIFO requires
$A_{e_{s+1}}\prec B_{o_s}$, while the other constraints give
\[
 B_{o_s}
 \mathrel{\prec_B} B_{e_1}
 \mathrel{\prec_H} A_{e_1}
 \mathrel{\prec_A} A_{e_{s+1}}
 \mathrel{\prec_{\rm LIFO}} B_{o_s}.
\]
The index ranges in the window table ensure that all events used in
these cycles exist. Thus a first divergence from $T$ is impossible.
Each boundary pair therefore has at most one complete H-valid merge.
Lemma~\ref{lem:per2-diagonal} covers every diagonal pair. Together these
cases cover all pairs $a=1$, $a=m$, or $a=b$. The constructions give
existence and D.5 gives full validity, completing the multiplicity-one
claim without using a count formula.

\subsection{D.10. Global-state-to-pivot bridge and
multiplicity}\label{d.10.-global-state-to-pivot-bridge-and-multiplicity}

Every globally valid \(MMMMVV\) state obeys all row-local constraints.
Proposition 6.1 therefore assigns its upper and lower restrictions
unique pivots \((a,b)\). This is a coverage/parameterization statement
only; it does not itself assert compatibility or existence.

Necessity shows the recovered pair is allowed. Sections D.3--D.5 give
one fully valid state for each allowed pair. Lemma~\ref{lem:per2-diagonal}
gives fixed-pair uniqueness on the diagonal, while Sections D.7--D.9
give fixed-pair uniqueness for the boundary families. These cases exhaust
$a=1$, $a=m$, or $a=b$. Hence the global-state-to-pair map is a
bijection and each compatible pair has multiplicity one. Only after this bridge may the compatible-pair
cardinality be identified with the state count.

\end{document}
